% 第6章 6.3.1：单位权星图 K_{1,3}（中心度3，叶度1）上两个单位脉冲。
% (a) 中心取1：能量=3，欧氏长度平方=1，度加权长度平方=d_c=3；
% (b) 一片叶取1：能量=1，欧氏长度平方=1，度加权长度平方=1。
\begin{tikzpicture}[x=1mm,y=1mm,font=\figurefont\footnotesize,text=FigureInk,
  vertex/.style={circle,draw=FigureStroke,fill=FigureBlueFill,line width=.8pt,minimum size=6mm,inner sep=0pt},
  hot/.style={vertex,draw=FigureYellow,fill=FigureYellowFill},
  edge/.style={draw=FigureStroke,line width=.8pt},
  note/.style={font=\FigureNoteFont\fontsize{7.5}{9}\selectfont,text=FigureMuted}]
\foreach \dx/\which/\tag in {0/0/(a) 中心取值为1,58/1/(b) 一片叶取值为1}{
 \begin{scope}[xshift=\dx mm]
  \ifnum\which=0
   \node[hot] (c) at (20,20) {$1$};
   \foreach \a [count=\i] in {90,210,330}{\node[vertex] (l\i) at ($(c)+(\a:11)$) {$0$};}
  \else
   \node[vertex] (c) at (20,20) {$0$};
   \node[hot] (l1) at ($(c)+(90:11)$) {$1$};
   \node[vertex] (l2) at ($(c)+(210:11)$) {$0$};
   \node[vertex] (l3) at ($(c)+(330:11)$) {$0$};
  \fi
  \foreach \i in {1,2,3}{\draw[edge] (c)--(l\i);}
  \node[anchor=west] at (-4,36) {\tag};
  \ifnum\which=0
   \node[anchor=west] at (-4,4.5) {能量$3$，$\lVert\bm g\rVert_2^2=1$，$\lVert\bm g\rVert_{\bm D}^2=3$};
  \else
   \node[anchor=west] at (-4,4.5) {能量$1$，$\lVert\bm g\rVert_2^2=1$，$\lVert\bm g\rVert_{\bm D}^2=1$};
  \fi
 \end{scope}}
\end{tikzpicture}
